\documentclass[11pt]{article}
% Unmodified official ACL style, embedded only for the standalone native editor.
% Upstream commit d5adc823ff0f80f98c80405ca0ab66c68e684409.
\begin{filecontents*}[overwrite]{acl.sty}
% This is the LaTex style file for *ACL.
% The official sources can be found at
%
%     https://github.com/acl-org/acl-style-files/
%
% This package is activated by adding
%
%    \usepackage{acl}
%
% to your LaTeX file. When submitting your paper for review, add the "review" option:
%
%    \usepackage[review]{acl}

\newif\ifacl@finalcopy
\newif\ifacl@anonymize
\newif\ifacl@linenumbers
\newif\ifacl@pagenumbers
\DeclareOption{final}{\acl@finalcopytrue\acl@anonymizefalse\acl@linenumbersfalse\acl@pagenumbersfalse}
\DeclareOption{review}{\acl@finalcopyfalse\acl@anonymizetrue\acl@linenumberstrue\acl@pagenumberstrue}
\DeclareOption{preprint}{\acl@finalcopytrue\acl@anonymizefalse\acl@linenumbersfalse\acl@pagenumberstrue}
\ExecuteOptions{final} % final copy is the default

% include hyperref, unless user specifies nohyperref option like this:
% \usepackage[nohyperref]{acl}
\newif\ifacl@hyperref
\DeclareOption{hyperref}{\acl@hyperreftrue}
\DeclareOption{nohyperref}{\acl@hyperreffalse}
\ExecuteOptions{hyperref} % default is to use hyperref
\ProcessOptions\relax

\typeout{Conference Style for ACL}

\usepackage{xcolor}

\ifacl@linenumbers
  % Add draft line numbering via the lineno package
  % https://texblog.org/2012/02/08/adding-line-numbers-to-documents/
  \usepackage[switch,mathlines]{lineno}

  % Line numbers in gray Helvetica 8pt
  \font\aclhv = phvb at 8pt
  \renewcommand\linenumberfont{\aclhv\color{lightgray}}

  % Zero-fill line numbers
  % NUMBER with left flushed zeros  \fillzeros[<WIDTH>]<NUMBER>
  \newcount\cv@tmpc@ \newcount\cv@tmpc
  \def\fillzeros[#1]#2{\cv@tmpc@=#2\relax\ifnum\cv@tmpc@<0\cv@tmpc@=-\cv@tmpc@\fi
    \cv@tmpc=1 %
    \loop\ifnum\cv@tmpc@<10 \else \divide\cv@tmpc@ by 10 \advance\cv@tmpc by 1 \fi
      \ifnum\cv@tmpc@=10\relax\cv@tmpc@=11\relax\fi \ifnum\cv@tmpc@>10 \repeat
    \ifnum#2<0\advance\cv@tmpc1\relax-\fi
    \loop\ifnum\cv@tmpc<#1\relax0\advance\cv@tmpc1\relax\fi \ifnum\cv@tmpc<#1 \repeat
    \cv@tmpc@=#2\relax\ifnum\cv@tmpc@<0\cv@tmpc@=-\cv@tmpc@\fi \relax\the\cv@tmpc@}%
  \renewcommand\thelinenumber{\fillzeros[3]{\arabic{linenumber}}}
  \AtBeginDocument{\linenumbers}

  \setlength{\linenumbersep}{1.6cm}

  % Bug: An equation with $$ ... $$ isn't numbered, nor is the previous line.

  % Patch amsmath commands so that the previous line and the equation itself
  % are numbered. Bug: multline has an extra line number.
  % https://tex.stackexchange.com/questions/461186/how-to-use-lineno-with-amsmath-align
  \usepackage{etoolbox} %% <- for \pretocmd, \apptocmd and \patchcmd

  \newcommand*\linenomathpatch[1]{%
    \expandafter\pretocmd\csname #1\endcsname {\linenomath}{}{}%
    \expandafter\pretocmd\csname #1*\endcsname {\linenomath}{}{}%
    \expandafter\apptocmd\csname end#1\endcsname {\endlinenomath}{}{}%
    \expandafter\apptocmd\csname end#1*\endcsname {\endlinenomath}{}{}%
  }
  \newcommand*\linenomathpatchAMS[1]{%
    \expandafter\pretocmd\csname #1\endcsname {\linenomathAMS}{}{}%
    \expandafter\pretocmd\csname #1*\endcsname {\linenomathAMS}{}{}%
    \expandafter\apptocmd\csname end#1\endcsname {\endlinenomath}{}{}%
    \expandafter\apptocmd\csname end#1*\endcsname {\endlinenomath}{}{}%
  }

  %% Definition of \linenomathAMS depends on whether the mathlines option is provided
  \expandafter\ifx\linenomath\linenomathWithnumbers
    \let\linenomathAMS\linenomathWithnumbers
    %% The following line gets rid of an extra line numbers at the bottom:
    \patchcmd\linenomathAMS{\advance\postdisplaypenalty\linenopenalty}{}{}{}
  \else
    \let\linenomathAMS\linenomathNonumbers
  \fi

  \AtBeginDocument{%
    \linenomathpatch{equation}%
    \linenomathpatchAMS{gather}%
    \linenomathpatchAMS{multline}%
    \linenomathpatchAMS{align}%
    \linenomathpatchAMS{alignat}%
    \linenomathpatchAMS{flalign}%
  }
\else
  % Hack to ignore these commands, which review mode puts into the .aux file.
  \newcommand{\@LN@col}[1]{}
  \newcommand{\@LN}[2]{}
  \newcommand{\nolinenumbers}{}
\fi

\PassOptionsToPackage{a4paper,margin=2.5cm,heightrounded=true}{geometry}
\RequirePackage{geometry}

\setlength\columnsep{0.6cm}
\newlength\titlebox
\setlength\titlebox{11\baselineskip}
% \titlebox should be a multiple of \baselineskip so that
% column height remaining fits an exact number of lines of text

\flushbottom \twocolumn \sloppy

% We're never going to need a table of contents, so just flush it to
% save space --- suggested by drstrip@sandia-2
\def\addcontentsline#1#2#3{}

\ifacl@pagenumbers
    \pagenumbering{arabic}
\else
    \thispagestyle{empty}
    \pagestyle{empty}
\fi

%% Title and Authors %%

\let\Thanks\thanks % \Thanks and \thanks used to be different, but keep this for backwards compatibility.

\newcommand\outauthor{%
    \begin{tabular}[t]{c}
    \ifacl@anonymize
        \bfseries Anonymous ACL submission
    \else
        \bfseries\@author
    \fi
    \end{tabular}}

% Mostly taken from deproc.
\AtBeginDocument{
\def\maketitle{\par
 \begingroup
   \def\thefootnote{\fnsymbol{footnote}}
   \twocolumn[\@maketitle]
   \@thanks
 \endgroup
 \setcounter{footnote}{0}
 \let\maketitle\relax
 \let\@maketitle\relax
 \gdef\@thanks{}\gdef\@author{}\gdef\@title{}\let\thanks\relax}
\def\@maketitle{\vbox to \titlebox{\hsize\textwidth
 \linewidth\hsize \vskip 0.125in minus 0.125in \centering
 {\Large\bfseries \@title \par} \vskip 0.2in plus 1fil minus 0.1in
 {\def\and{\unskip\enspace{\rmfamily and}\enspace}%
  \def\And{\end{tabular}\hss \egroup \hskip 1in plus 2fil
           \hbox to 0pt\bgroup\hss \begin{tabular}[t]{c}\bfseries}%
  \def\AND{\end{tabular}\hss\egroup \hfil\hfil\egroup
          \vskip 0.25in plus 1fil minus 0.125in
           \hbox to \linewidth\bgroup\large \hfil\hfil
             \hbox to 0pt\bgroup\hss \begin{tabular}[t]{c}\bfseries}
  \hbox to \linewidth\bgroup\large \hfil\hfil
    \hbox to 0pt\bgroup\hss
  \outauthor
   \hss\egroup
    \hfil\hfil\egroup}
  \vskip 0.3in plus 2fil minus 0.1in
}}
}

% margins and font size for abstract
\renewenvironment{abstract}%
  {\begin{center}\large\textbf{\abstractname}\end{center}%
    \begin{list}{}%
      {\setlength{\rightmargin}{0.6cm}%
        \setlength{\leftmargin}{0.6cm}}%
      \item[]\ignorespaces%
      \@setsize\normalsize{12pt}\xpt\@xpt
  }%
  {\unskip\end{list}}

% Resizing figure and table captions - SL
% Support for interacting with the caption, subfigure, and subcaption packages - SL
\RequirePackage{caption}
\DeclareCaptionFont{10pt}{\fontsize{10pt}{12pt}\selectfont}
\captionsetup{font=10pt}

\RequirePackage{natbib}
% for citation commands in the .tex, authors can use:
% \citep, \citet, and \citeyearpar for compatibility with natbib, or
% \cite, \newcite, and \shortcite for compatibility with older ACL .sty files
\renewcommand\cite{\citep}  % to get "(Author Year)" with natbib
\newcommand\shortcite{\citeyearpar}% to get "(Year)" with natbib
\newcommand\newcite{\citet} % to get "Author (Year)" with natbib
\newcommand{\citeposs}[1]{\citeauthor{#1}'s (\citeyear{#1})} % to get "Author's (Year)"

\bibliographystyle{acl_natbib}

% Bibliography

% Don't put a label in the bibliography at all.  Just use the unlabeled format
% instead.
\def\thebibliography#1{\vskip\parskip%
\vskip\baselineskip%
\def\baselinestretch{1}%
\ifx\@currsize\normalsize\@normalsize\else\@currsize\fi%
\vskip-\parskip%
\vskip-\baselineskip%
\section*{References\@mkboth
 {References}{References}}\list
 {}{\setlength{\labelwidth}{0pt}\setlength{\leftmargin}{\parindent}
 \setlength{\itemindent}{-\parindent}}
 \def\newblock{\hskip .11em plus .33em minus -.07em}
 \sloppy\clubpenalty4000\widowpenalty4000
 \sfcode`\.=1000\relax}
\let\endthebibliography=\endlist


% Allow for a bibliography of sources of attested examples
\def\thesourcebibliography#1{\vskip\parskip%
\vskip\baselineskip%
\def\baselinestretch{1}%
\ifx\@currsize\normalsize\@normalsize\else\@currsize\fi%
\vskip-\parskip%
\vskip-\baselineskip%
\section*{Sources of Attested Examples\@mkboth
 {Sources of Attested Examples}{Sources of Attested Examples}}\list
 {}{\setlength{\labelwidth}{0pt}\setlength{\leftmargin}{\parindent}
 \setlength{\itemindent}{-\parindent}}
 \def\newblock{\hskip .11em plus .33em minus -.07em}
 \sloppy\clubpenalty4000\widowpenalty4000
 \sfcode`\.=1000\relax}
\let\endthesourcebibliography=\endlist

% sections with less space
\def\section{\@startsection {section}{1}{\z@}{-2.0ex plus
    -0.5ex minus -.2ex}{1.5ex plus 0.3ex minus .2ex}{\large\bfseries\raggedright}}
\def\subsection{\@startsection{subsection}{2}{\z@}{-1.8ex plus
    -0.5ex minus -.2ex}{0.8ex plus .2ex}{\normalsize\bfseries\raggedright}}
%% changed by KO to - values to get the initial parindent right
\def\subsubsection{\@startsection{subsubsection}{3}{\z@}{-1.5ex plus
   -0.5ex minus -.2ex}{0.5ex plus .2ex}{\normalsize\bfseries\raggedright}}
\def\paragraph{\@startsection{paragraph}{4}{\z@}{1.5ex plus
   0.5ex minus .2ex}{-1em}{\normalsize\bfseries}}
\def\subparagraph{\@startsection{subparagraph}{5}{\parindent}{1.5ex plus
   0.5ex minus .2ex}{-1em}{\normalsize\bfseries}}

% Footnotes
\footnotesep 6.65pt %
\skip\footins 9pt plus 4pt minus 2pt
\def\footnoterule{\kern-3pt \hrule width 5pc \kern 2.6pt }
\setcounter{footnote}{0}

% Lists and paragraphs
\parindent 1em
\topsep 4pt plus 1pt minus 2pt
\partopsep 1pt plus 0.5pt minus 0.5pt
\itemsep 2pt plus 1pt minus 0.5pt
\parsep 2pt plus 1pt minus 0.5pt

\leftmargin 2em \leftmargini\leftmargin \leftmarginii 2em
\leftmarginiii 1.5em \leftmarginiv 1.0em \leftmarginv .5em \leftmarginvi .5em
\labelwidth\leftmargini\advance\labelwidth-\labelsep \labelsep 5pt

\def\@listi{\leftmargin\leftmargini}
\def\@listii{\leftmargin\leftmarginii
   \labelwidth\leftmarginii\advance\labelwidth-\labelsep
   \topsep 2pt plus 1pt minus 0.5pt
   \parsep 1pt plus 0.5pt minus 0.5pt
   \itemsep \parsep}
\def\@listiii{\leftmargin\leftmarginiii
    \labelwidth\leftmarginiii\advance\labelwidth-\labelsep
    \topsep 1pt plus 0.5pt minus 0.5pt
    \parsep \z@ \partopsep 0.5pt plus 0pt minus 0.5pt
    \itemsep \topsep}
\def\@listiv{\leftmargin\leftmarginiv
     \labelwidth\leftmarginiv\advance\labelwidth-\labelsep}
\def\@listv{\leftmargin\leftmarginv
     \labelwidth\leftmarginv\advance\labelwidth-\labelsep}
\def\@listvi{\leftmargin\leftmarginvi
     \labelwidth\leftmarginvi\advance\labelwidth-\labelsep}

\abovedisplayskip 7pt plus2pt minus5pt%
\belowdisplayskip \abovedisplayskip
\abovedisplayshortskip  0pt plus3pt%
\belowdisplayshortskip  4pt plus3pt minus3pt%

% Less leading in most fonts (due to the narrow columns)
% The choices were between 1-pt and 1.5-pt leading
\def\@normalsize{\@setsize\normalsize{11pt}\xpt\@xpt}
\def\small{\@setsize\small{10pt}\ixpt\@ixpt}
\def\footnotesize{\@setsize\footnotesize{10pt}\ixpt\@ixpt}
\def\scriptsize{\@setsize\scriptsize{8pt}\viipt\@viipt}
\def\tiny{\@setsize\tiny{7pt}\vipt\@vipt}
\def\large{\@setsize\large{14pt}\xiipt\@xiipt}
\def\Large{\@setsize\Large{16pt}\xivpt\@xivpt}
\def\LARGE{\@setsize\LARGE{20pt}\xviipt\@xviipt}
\def\huge{\@setsize\huge{23pt}\xxpt\@xxpt}
\def\Huge{\@setsize\Huge{28pt}\xxvpt\@xxvpt}

% The hyperref manual (section 9) says hyperref should be loaded after natbib
\ifacl@hyperref
  \PassOptionsToPackage{breaklinks}{hyperref}
  \RequirePackage{hyperref}
  % make links dark blue
  \definecolor{darkblue}{rgb}{0, 0, 0.5}
  \hypersetup{colorlinks=true, citecolor=darkblue, linkcolor=darkblue, urlcolor=darkblue}
\else
  % This definition is used if the hyperref package is not loaded.
  % It provides a backup, no-op definiton of \href.
  % This is necessary because \href command is used in the acl_natbib.bst file.
  \def\href#1#2{{#2}}
  \usepackage{url}
\fi
\end{filecontents*}
\usepackage[review]{acl}
\usepackage{times,latexsym}
\usepackage[T1]{fontenc}
\usepackage[utf8]{inputenc}
\usepackage{microtype}
\usepackage{amsmath,amssymb,amsthm,booktabs}
\newtheorem{theorem}{Theorem}
\newtheorem{lemma}{Lemma}
\newtheorem{proposition}{Proposition}
\newcommand{\R}{\mathbb{R}}
\newcommand{\Q}{\mathbb{Q}}
\newcommand{\calD}{\mathcal{D}}
\newcommand{\calS}{\mathcal{S}}
\newcommand{\Qfix}{Q_{\mathrm{fix}}}
\newcommand{\evidence}[1]{(supplementary evidence #1)}
\title{Certified Calibration Repair from Compressed Fixed-Anchor Features}
\author{Anonymous submission}
\date{}
\begin{document}
\maketitle

\begin{abstract}
Can a smaller calibration archive still recover the exact quantized model after
records are removed? We study a fixed-teacher reconstruction target with exact
dyadic feature interpretation, fixed normalization and grids, and deterministic
adaptive rounding. Source-local compressed enclosures support universal decision
certificates, bounded exact replay, and canonical successor states. We prove
conditional code and state correctness; resource-bounded completion and practical
benefit require separate evidence.
On DistilGPT2 with $1664\to1536\to768$ calibration tokens, five service
arms on a disjoint development root match fresh reconstruction for all
42.47 million codes and each canonical successor state. With the original
checkpoint as teacher, compression saves 9.24\% and 8.70\% of live state against
a Zstandard-augmented lossless cache. A complete optimization screen preserves
exact outputs: compressed requests are 0.90\% faster at the first endpoint and
12.73\% slower at the second than the equally optimized cache, with one
observation per cell. A hybrid exact-Gram control is faster and smaller at the first endpoint in the original comparison,
while requiring source access. On 320 outcome-blind selected articles, the
fixed-teacher target has 0.54\% and 0.71\% higher pooled perplexity than matched
sequential calibration; both endpoints pass the frozen 5\% pooled and 20\%
maximum-article tolerances, including paired bootstrap upper bounds.
The earlier nearest-anchor target failed with a 5.92\% perplexity regression beyond a 5\% tolerance.
These results establish complete-state correctness and a modest storage
tradeoff with a small measured quality cost. General acceleration and
quality transfer beyond this bounded frame remain unestablished.
\end{abstract}

\section{Introduction}
Post-training quantization uses calibration data to select low-bit model weights.
When calibration membership changes, an application may need a replacement that
matches an explicitly defined retained-data reconstruction. Repeating the entire
calibration computation is one option. Keeping all exact feature values is
another. A smaller archive is useful only if it preserves enough information to
recover the required discrete decisions, or can detect when additional source
access is needed.

Our question is therefore about the joint cost of \emph{preparing, retaining,
accessing, and updating calibration evidence}. A repair method is not useful
merely because it is faster than an unoptimized cold constructor. An equally
informed indexed constructor can execute the same repair procedure. Strong
comparisons must include exact cached features, exact additive statistics, and
fresh reconstruction under matching output and access contracts.

We examine compressed, source-local feature enclosures. Compression discards
exact feature words from the live descriptor while retaining intervals that
contain them. A numerical proposal suggests coefficients or codes; a separate
outward verifier must certify each decision for every realization of the feature
box. An unresolved stage can trigger an admitted exact replay and point
certificate. If those checks do not complete, the service refuses rather than
returning an approximate replacement.

The counterfactual is deliberately restricted. A teacher fixed from the pretrained
checkpoint alone, either nearest-grid or original-checkpoint, produces every
calibration feature. Repaired weights do not change later calibration features. This differs
from ordinary sequential calibration and from stock floating-point GPTQ
implementations~\cite{gptq}. It also concerns calibration membership, not removal
of information from the checkpoint's pretraining.

We specify an adaptive-code recovery contract and a canonical successor state,
then test their complete storage, access and latency costs. Fixed features,
sufficient statistics and decision stability are established ingredients. Our
contribution is their use in a checked repair service with uncertain archived
features, together with an explicit account of its limits. On a disjoint
calibration root, all five service arms match fresh reconstruction; compression
saves live storage but initially takes longer than the strongest lossless cache.
A complete optimization screen narrows that gap without a uniform latency win.
Untouched quality measurements show a small perplexity cost within the frozen tolerances.
The earlier failed quality guard and adaptive teacher choice remain part of the
evidence~\evidence{E03, E12, E13, E14, E16, E17}.
\section{An explicit reconstruction target}
\label{sec:target}
Let $W$ be a fixed pretrained checkpoint and let $\calD_0$ be an original set of
identified token records. Fix the tokenizer, record boundaries, stage order,
coordinate order, ridge parameters, row grids, and finite feature evaluator.
The original normalization $N_0>0$ is part of this specification; deletion does
not replace it with the number of surviving tokens.

Fix a checkpoint-only teacher $W_{\mathrm{teacher}}=T(W)$ before selecting
retained membership. The nearest-grid instantiation rounds each stage's base
weights on its frozen grid with lower-code midpoint ties; the original-checkpoint
instantiation uses $T(W)=W$. These define distinct targets, whose empirical
results cannot be transferred between them. For source $j$, let
\[
 X_{\ell j}=F_\ell(W_{\mathrm{teacher}},j)
       \in\Q^{d_\ell\times t_j}.
\]
Here $F_\ell$ is the declared finite CPU computation, followed by exact
interpretation of each successful binary64 output word as a dyadic rational.
It is not an ideal real-arithmetic transformer. Unsupported or unresolved finite
operations cause refusal, so this evaluator is partial. Each source is evaluated
independently; its features cannot depend on another source's membership.

For a retained set $R\subseteq\calD_0$, concatenate its features in canonical
source order to obtain $X_\ell(R)$, and define
\begin{equation}
 H_\ell(R)=\lambda_\ell I+
   \frac{X_\ell(R)X_\ell(R)^\top}{N_0},
 \qquad \lambda_\ell>0.
 \label{eq:metric}
\end{equation}
A target with stage-specific original normalizations can be defined similarly;
all empirical configurations discussed here use one frozen token count.

\paragraph{Exact adaptive rounding.}
For one stage, write its rational positive-definite metric as a reverse LDL
factorization
\[
 H=L^\top\operatorname{diag}(\tau)L,
 \qquad B=I-L,
\]
where $L$ is unit lower triangular and every pivot $\tau_i$ is positive. For each
base-weight row $w$, process coordinates in their frozen order:
\begin{align}
 v_i&=w_i-\sum_{h<i} B_{ih}(w_h-q_h),\label{eq:recurrence}\\
 q_i&=\operatorname{nearest}_{G_i}(v_i).
 \label{eq:rounding}
\end{align}
Stage and output-row indices are suppressed in this notation. Each output row
uses the smallest checkpoint-only dyadic scale $s_{\ell,r}$ with at most 24
significant binary digits whose grid
$G_{\ell,r}=s_{\ell,r}\{-8,\ldots,7\}$ covers its base weights and has exactly
representable finite codes and midpoints. The scale is shared across coordinates.
Each finite increasing rational grid $G_i$ is fixed before deletion. Nearest
rounding saturates at endpoints; an exact midpoint selects the lower code. The
collection of all stage outputs defines $\Qfix(R)$ wherever the required feature
evaluations succeed. Positive ridge makes the mathematical suffix systems
invertible, but does not guarantee that a bounded numerical implementation will
certify every decision.

\paragraph{What remains fixed.}
The checkpoint, teacher rule, original normalization, grids, arithmetic target, and
ordering belong to target identity. Changing any of them requires a new target.
Fresh reconstruction used as an oracle must retain the same original
normalization. Matching a separately retuned or sequentially recalibrated model
would answer a different question.

\section{Repair from compressed evidence}
\label{sec:method}
Trusted preparation evaluates each original source once through all stages.
For each $X_{\ell j}$ it constructs a canonical descriptor $D_{\ell j}$ whose
public decoding is an entrywise box
\begin{equation}
 \underline X_{\ell j}\leq X_{\ell j}\leq\overline X_{\ell j}.
 \label{eq:containment}
\end{equation}
Preparation checks this containment against the actual finite values. Descriptor
metadata binds the source ID, token identity, stage, target, codec settings, and
source-word commitment. A hash detects changed bytes; it does not independently
prove containment or authenticate an untrusted execution host.

\paragraph{Stored enclosures.}
The implemented codec traverses each source-stage matrix in token-major order,
grouping 256 values per block. Each block uses a shared power-of-two step and
signed 48-bit integer cell indices. Two-bit tags distinguish closed cells, exact
grid values, and exact binary64 escapes near overflow. Decoding constructs exactly
representable dyadic endpoints, and preparation checks every original value
against them. These 48 bits describe calibration evidence; deployed weight codes
use four bits. Metadata and escapes also contribute to the reported state size.

A request validates the predecessor, retained membership, deletion IDs, and all
target and provenance bindings. Retained source-local descriptors remain
unchanged. At each stage the service presents their decoded boxes to a universal
certificate. The numerical solver may use floating-point proposals, but proposals
alone never establish correctness. Directed residual bounds and decision-cell
checks supply the evidence.

\paragraph{Decision cells.}
For $G_i=\{g_0<\cdots<g_{k-1}\}$, let
$m_a=(g_a+g_{a+1})/2$. The cell selecting an interior code $g_a$ is
$(m_{a-1},m_a]$; endpoint cells extend to infinity. Thus a closed enclosure
$[\underline v_i,\overline v_i]$ certifies $g_a$ only when it lies wholly within
that cell. A lower boundary is strict and an upper boundary is inclusive. This
asymmetry is required by the lower-code tie rule, including exact midpoint cases.
Previously certified codes are used when bounding the next adaptive decision.

\paragraph{Bounded fallback.}
If a box certificate refuses, the service may regenerate the retained sources'
exact finite features and invoke an admitted point certificate. Every traversed
ancestor is charged, including stages whose compressed certificates already
accepted. Replayed source words and their containment must match trusted
preparation. Repeated stage failures may reuse an already regenerated source
within the same request, with that access made explicit. Temporary exact features
do not silently replace the canonical compressed representation. Exhausted CPU,
memory, refinement, or replay allowances cause refusal without a complete model
release.

\subsection{Stage-specific numerical representations}
For width $d$, retained token count $T$, and output-row count $r$, token-space
coefficient construction has a leading structural term proportional to $dT^2$.
Feature-space construction has terms proportional to $d^2T+d^3$, followed by
row work proportional to $rd^2$. These expressions describe arithmetic schedules,
not measured runtime. Interval width, repeated verification, decoding, source
replay, and native implementation constants can reverse finite-size rankings.

The feature-space route builds directed Gram enclosures from actual feature
boxes, or accepts an exact Gram with trusted source lineage. It must not infer
positive semidefiniteness for every arbitrary matrix in an entrywise Gram box.
Every \emph{realized feature Gram} is positive semidefinite; that is the premise
used by the certificate.

For $K=XX^\top$, $\beta=\lambda N_0$, and a suffix $S$, let
$A=\beta I+K_{S,S}$. If a proposal $\widehat z$ has a verified residual bound
$\|b-A\widehat z\|_2\leq\eta$, then
\begin{equation}
 \|A^{-1}b-\widehat z\|_2\leq\frac{\eta}{\beta}.
 \label{eq:residual}
\end{equation}
This follows from $A\succeq\beta I$. For uncertain features the residual bound
must hold uniformly over their box, including directed product and summation
errors. Such bounds support coefficient enclosures and then the adaptive cell
checks; they do not make an unverified nominal solve exact.

Our dispatcher admits the complete bounded solver schedule and its
explicit arrays before feature access. It prefers token space when $4T\leq d$
and streamed primal otherwise, subject to admission. This is a declared
engineering rule, not a learned optimal crossover. Compatible controls receive
the same kernels and choices. A trusted native exact Gram uses the same
feature-space certificate; its integer accumulator and archive parsing have
separate bounds.

\section{Conditional correctness and state}
The following statements concern the specified target and access contract. They
are not claims of universal completion, speed, cryptographic authenticity, or
novelty of interval verification.

\begin{theorem}[Conditional code correctness]
\label{thm:codes}
Assume that retained features are defined by the frozen source-local evaluator,
trusted descriptors satisfy~\eqref{eq:containment}, and every accepted rounding
certificate is valid for all feature realizations in its supplied box, conditional
on the accepted prefix. Assume that any fallback reproduces those same finite
features and certifies~\eqref{eq:recurrence}--\eqref{eq:rounding}. If the service
successfully commits all required stages, its model is $\Qfix(R)$.
\end{theorem}
\begin{proof}
Source locality makes each retained source's feature matrix equal to its
retained-only reconstruction. These matrices lie inside the prepared boxes.
For the first coordinate, a universally valid cell certificate therefore selects
the exact target code. Suppose the preceding codes equal the target prefix. The
certificate's conditional premise is then satisfied, so its next accepted code
also equals the target. Induction establishes every row. A successful fallback
uses the same features and exact point target, hence has the same conclusion.
The fixed teacher is independent of calibrated outputs, so this argument applies
to every stage without changing a later stage's defining inputs. The complete
extent check establishes equality of the entire committed model.
\end{proof}

\paragraph{Canonical service state.}
Fix the codec, preparer, provenance policy, source ordering, target serialization,
and representation rule. Let $E(R)$ contain either retained canonical source
descriptors or the canonical exact pooled Grams required by that rule. Define
\begin{equation}
 \calS(R)=\operatorname{encode}\bigl(
   \text{target},\Qfix(R),R,E(R),\text{provenance}\bigr).
 \label{eq:state}
\end{equation}
Required tokens or source-retrieval obligations and external trust sidecars are
part of the service's storage/access accounting. The common pretrained checkpoint
and runtime dependencies remain required even if reported separately.
State provenance means retained-record provenance and target/codec bindings.
Execution logs and execution-provenance artifacts remain external; canonical
state equality alone does not compare them or the external Gram-trust sidecars.
The packed model artifact contains code indices for all 24 calibrated affine
stages; unchanged checkpoint components are still needed for inference.

\begin{lemma}[Exact source subtraction]
For fixed source-local features, $K(R)=\sum_{j\in R}X_jX_j^\top$. Exact addition and
subtraction of trusted source contributions, with original normalization
unchanged, produce the same rational Gram as fresh retained-source accumulation.
\end{lemma}
\begin{proof}
The summands are unchanged by other sources' deletion. Finite sums over disjoint
source sets are additive in exact rational arithmetic. Subtracting precisely the
removed summands therefore leaves exactly the retained sum. Canonical integer
scaling, source ordering, and serialization remove representation dependence on
addition or deletion order.
\end{proof}

\begin{theorem}[Canonical successor equality]
\label{thm:state}
Under Theorem~\ref{thm:codes}, suppose source descriptors are deterministic and
retained unchanged, exact Gram updates satisfy the preceding lemma, and all
compared executions use the same policies in~\eqref{eq:state}. A successful repair
returns the abstract canonical model and state $\calS(R)$, and therefore the same
bytes as any successful certified fresh retained-input reconstruction in the
same representation. This does not guarantee that a bounded fresh numerical
reconstruction completes. The equality conclusion holds after any
sequence of successful valid deletions.
\end{theorem}
\begin{proof}
The model agrees by Theorem~\ref{thm:codes}. Each retained descriptor agrees because
its source-local inputs and deterministic encoder agree. Each pooled Gram agrees
by exact subtraction and canonicalization. Membership, target, provenance, and
serialization policies are identical, so every serialized component agrees.
Applying this argument to each valid successor proves the sequence claim.
\end{proof}

State equality is representation-specific: a Gram state need not equal a
compressed-descriptor state. Changing the preparer or provenance format can also
change state bytes while leaving model codes unchanged. Implementation tests must
therefore compare independently reconstructed actual payloads, not only output
hashes or parser roundtrips. The theorem assumes precisely the premises those
checks are intended to exercise.

\section{Complete cost and access accounting}
A repair may reduce live storage while increasing request latency. Let $P_a$ be
initial preparation for representation $a$, and $C_{a,k}$ its complete cost on the
$k$th actual successor transition. For a declared initial-model convention,
\begin{equation}
 L_a(m)=P_a+\sum_{k=1}^{m}C_{a,k}
 \label{eq:lifetime}
\end{equation}
is its preparation-inclusive lifetime cost. A break-even observation requires
comparing these complete sums under matched access and output contracts. Repeating
alternative deletions from one unchanged root is not a successor-state sequence.
No theorem above bounds~\eqref{eq:lifetime} by a competitor's cost.

Request clocks must account for loading and validating the predecessor, source
access, exact accumulation, every certificate attempt, fallback, serialization,
complete output, and disposal. Shared checkpoint load, native compilation,
controller work, and independent audits need explicit placement rather than an
implicit zero. Nested timings cannot be added to their enclosing timer as if they
were disjoint. Refused runs retain their full observed cost and are not successful
speed observations.

An exact pooled Gram can be compact at large token counts, but deleting a record
requires that record's contribution. If individual contributions are not stored,
the service must retain access to deleted-source tokens or another way to
regenerate them. The pooled Gram and a membership digest alone do not supply that
information. Compressed fallback similarly requires retained-source access.

Storage includes model codes, descriptors or Grams, targets, tokens or retrieval
indexes, provenance, and required trust metadata. We report the shared checkpoint
separately and in a complete deployment total. An exported model already embedded
in a state must not be counted twice. Audit archives are separate physical data;
logical deletion from live calibration state does not erase those archives.

\section{Development history}
\label{sec:history}
Historical pilots include WikiText-2~\cite{wikitext}.
Earlier stage-scale and tiny complete-model pilots established integration and
actual successor equality, while exposing strong cache and Gram controls.
They used development inputs and do not supply independent confirmation.
The reproduction record retains their complete tables, runtime identities,
certificate refusals, replay costs and failed attempts~\evidence{E01, E02, E04, E05, E06}.
All sixty historical quality-article exclusions remain protected. Same-model
GPU likelihood agreement validates that evaluator comparison; it does not
establish exact GPU calibration or a speed benefit.

\section{Related work and contribution boundary}
GPTQ develops post-training quantization using approximate second-order
information~\cite{gptq}. Our reconstruction contract must be stated independently:
finite source-local fixed-anchor features, original normalization, exact dyadic
interpretation, frozen grids, and lower midpoint ties. Agreement with this target
does not certify agreement with an unspecified GPTQ implementation or with a
sequential calibration target whose downstream activations change.

Exact sufficient statistics, ridge systems, low-rank identities, and margin
stability are established methods. QSS~\cite{qss} freezes a learned schema and
maintains additive prediction corrections over quantized regions, distinguishing
label deletion from instance deletion and accounting for schema rebuilds. Our
target instead reconstructs adaptive weight codes under a fixed calibration
constructor. MEDU~\cite{medu} compresses stored distributed-training updates and
bounds deviation from retraining. Our conditional contract concerns exact discrete
recovery from compressed feature enclosures, not the first use of lossy storage
for unlearning.

\paragraph{Quantization-enabled exact deletion.}
Ginart et al.~\cite{ginart} retain metadata for quantized $k$-means, check whether
deletion changes quantized centroids, update metadata when stable, and otherwise
retrain. Their randomized-lattice analysis concerns expected deletion costs.
Thus stability checks, exact deletion aided by quantization, retraining fallback,
and unchanged models with updated metadata are established prior work.

Exact-Fun~\cite{exactfun} changes federated training to a quantized procedure,
retains historical models, corrects deleted gradients, checks recomputed
quantized parameters, and retrains from an affected round. Its full author-version
Theorem 2 proof uses uniform within-cell positions and independent coordinates;
the overlap calculation also uses a uniform perturbation model. Its probability
formula is not a distribution-free consequence of a displacement bound alone.
Our intended distinction is universal recovery of the specified adaptive codes
over uncertain feature evidence, including changed codes, rather than the
general fact that rounding can leave an output unchanged after deletion.

\paragraph{Executable certificates.}
ExecCert~\cite{execcert} separates candidate generation, the release contract, and
the verifier. It studies executable residual certificates, transfer to exact
references, guarded updates and rebuilds, decision agreement, and lifecycle costs.
Its Section 4.4 explicitly discusses finite output formats and codebooks. Our
adaptive weight-code contract over uncertain calibration features specializes a
different obligation from ridge-relative or fixed-readout tolerances; this does
not establish that the broader framework cannot express that contract.

QEBVerif~\cite{qebverif} combines differential reachability bounds with MILP
fallback to verify inference-output differences between a network and its quantized
counterpart. Our property concerns the codes produced by a calibration algorithm
when its feature evidence is uncertain.

These comparisons restrict our contribution to a specified adaptive-code recovery
contract and canonical successor state. They do not establish a new general
verification principle. We do not claim empirical superiority over these external
systems or reproduce their implementations.

\section{Complete-model evidence}
We separate exact reconstruction, prediction quality and complete service cost.
We use Hugging Face's DistilGPT2 (approximately 82 million total parameters)
and English C4 web text~\cite{distilgpt2,c4}. The development workload uses
thirteen training records of 128 tokens each, then deletes one record followed
by six. Original normalization stays 1664.
One CPU thread evaluates all 24 stages and 42,467,328 model codes. Each fresh
oracle independently traverses retained tokens but shares the numerical solver;
it is not a separate implementation replication.

\paragraph{Preserved failures and adaptive choices.}
The first larger-root preparation exceeded its 4,800-CPU-second cap after
11 of 24 certified stages, charging 4,801 seconds and cancelling seven
dependents~\evidence{E07}. No complete model resulted. A later wide-stage
gate verified optimized and historical point-solver byte equality, but its
40-bit compressed arm refused after 79 rows. A separately registered 48-bit
follow-up certified the same stage; subsequent service experiments use that
adaptively selected precision~\evidence{E08, E09}. These are engineering decisions,
not prospective confirmation. The failed registration, every refusal and their
costs remain in the record~\evidence{E10, E12}.

\paragraph{Nearest-anchor service and failed quality guard.}
The nearest-rounded teacher's complete service matches fresh model/state bytes
at all eight arm--endpoint comparisons, but its quality guard fails~\evidence{E12}.
Eight previously exposed articles supply 1,016 predictions per model~\evidence{E11, E12}.
The nearest-anchor model's PPL is 55.14666, versus 52.06342 for sequential calibration:
ratio 1.059221 exceeds the unchanged 1.05 guard.
Loss is higher on all eight articles. The maximum article ratio, 1.174280,
passes the separate 1.20 guard, but the overall gate fails. All service arms
return this same model, so the regression belongs to the target. Separate
CPU/CUDA likelihood agreement preserves this adverse finding.

\paragraph{Adaptive full-precision-anchor screen (V59).}
Following that failure, a separately registered candidate replaced the
nearest-rounded teacher with the original checkpoint~\evidence{E13}. The deployed
weights still use four-bit codes. Retained tokens, normalization, grids and the
eight exposed articles stayed fixed. Its pooled PPL is 51.85008 (mean NLL 3.948356),
giving candidate/sequential ratio 0.995902 and maximum article ratio 1.065598.
Both unchanged guards, 1.05 pooled and 1.20 per article, pass.
Candidate PPL remains 13.71\% above the unquantized checkpoint. This adaptive
screen motivated the later service; it is not untouched evidence or a quality-gain
claim. Construction timings and article-level results remain in E13.

\paragraph{Original-checkpoint service (V61; recovered audits V62).}
All eight arm--endpoint comparisons match fresh model and canonical-state bytes;
the final-six models equal the screened candidate~\evidence{E14}.
Compression saves 9.41\% and 8.86\% of live state against cached lossless,
but requests take 38.62\% and 41.74\% longer. The original primary audit timed out
and its dependent was cancelled. Separate artifact-only recovery used the
unchanged validators without repeating numerical work. The failed original audit
consumed another 1,811.125 CPU seconds. Full preparation, request, deployment and
recovery costs remain in E14. These measurements concern the earlier calibration
root; the disjoint-root service and untouched quality below are separate runs.

\paragraph{Stronger lossless artifact control (V63).}
The stronger cache adds plain and byte-shuffled Zstandard3 to raw, zlib6 and
shuffled-zlib6 encodings, selecting the smallest complete descriptor with fixed
ties. It charges every candidate encoder and validation. An artifact-only audit
restored the original state bytes for all three memberships~\evidence{E15}.
This control copies existing model codes and supplies no repair latency result.
Its complete service is evaluated on the disjoint root below.

\paragraph{Disjoint-root complete service (V67).}
A separately selected root uses thirteen previously unexposed training records
and the same deletion rule, teacher, grids, solver and 48-bit archive precision.
All ten arm--endpoint model/state comparisons match fresh retained-input
reconstruction, including the complete modern-lossless service~\evidence{E16}.
Table~\ref{tab:v67} reports the complete requests, live state and source access.
Against cached lossless, compression saves 9.38\% and 8.83\% of live state but
requests take 17.82\% and 46.46\% longer. Against the modern cache, savings are
9.24\% and 8.70\%, with requests 13.69\% and 50.48\% longer; checkpoint-inclusive
savings fall to 4.71\% and 3.10\%. The added codec installation occupies
23,870,575 bytes, accounted separately from these checkpoint-inclusive totals.
Hybrid remains faster and smaller at the first endpoint but reads source tokens;
compression requires no source traversal on this trajectory. At the second
endpoint it is faster and smaller than hybrid, while both caches remain faster.
No stage refusal or source-replay fallback occurs. Compression uses bounded
numerical fallbacks on 68 and 2 rows, which are distinct from source replay.
The numerical workers charge 47,767 CPU seconds; all fifteen scheduler
allocations, including both audits, consume 51,356.452 CPU seconds.
These clocks overlap and must not be added. The complete comparison retains one
fixed order per endpoint, so it provides neither timing uncertainty nor a
population-level reliability estimate. The untouched quality comparison follows below.

\begin{table*}[t]
\centering\small
\begin{tabular}{rlrrrr}
\toprule
Records & Service arm & Request (s) & Live (MB) & + checkpoint (MB) & Access \\
\midrule
12 & Cached lossless & 4,419.007 & 368.003 & 720.828 & 0 \\
 & Modern lossless & 4,579.496 & 367.456 & 720.282 & 0 \\
 & Cold reconstruction & 5,973.008 & 368.003 & 720.828 & 288 \\
 & Hybrid exact Gram & 4,921.935 & 318.918 & 671.743 & 24 \\
 & Compressed 48-bit & 5,206.569 & 333.497 & 686.322 & 0 \\
\midrule
6 & Cached lossless & 922.484 & 195.615 & 548.441 & 0 \\
 & Modern lossless & 897.841 & 195.330 & 548.155 & 0 \\
 & Cold reconstruction & 1,619.473 & 195.615 & 548.441 & 144 \\
 & Hybrid exact Gram & 1,694.311 & 220.247 & 573.072 & 144 \\
 & Compressed 48-bit & 1,351.098 & 178.333 & 531.158 & 0 \\
\bottomrule
\end{tabular}
\caption{Disjoint-root complete service (V67). Decimal MB include trust sidecars;
the next column adds the shared checkpoint, excluding runtime/codec installation.
Access counts source--stage record traversals during the request. Zero means no
source traversal, not zero computation. Each row is one fixed-order observation;
there is no timing confidence interval. Full preparation and transaction costs
are in E16.}
\label{tab:v67}
\end{table*}

\paragraph{Complete optimization screen (V89).}
We next keep the same thirteen-record root, target, two deletions and four-cell
comparison. Both optimized arms receive the compatible token-neighbor kernel;
compressed storage additionally receives a native payload decoder, while the
modern cache reuses decoded canonical descriptors within a bounded request cache.
The two changes in each arm form a package, so this screen does not isolate a
decoder or cache effect. Every request starts in a fresh process. The order is
fixed prospectively and reversed at the second endpoint; OS cache state remains
uncontrolled. An earlier execution failed at the storage quota. A separately
registered recovery preserves that failure and its unresolved resource debit.

All eight complete state/trust/model comparisons match independent fresh
reconstruction, with no source traversal or replay fallback. Table~\ref{tab:v89}
uses complete host time, including admission, container startup and cleanup.
The compressed package reduces its observed request times by 11.53\% and 36.17\%;
the modern package reduces them by 0.69\% and 19.93\%. Optimized compression is
0.90\% faster than optimized modern at twelve records but 12.73\% slower at six.
Its live and checkpoint-inclusive storage sizes remain those in Table~\ref{tab:v67}.
These single observations do not establish stable latency or general speedup.

\begin{table}[t]
\centering\small
\begin{tabular}{rlrr}
\toprule
Records & Representation & Reference (s) & Optimized (s) \\
\midrule
12 & Compressed & 4,444.116 & 3,931.646 \\
 & Modern lossless & 3,995.241 & 3,967.523 \\
6 & Compressed & 1,328.624 & 848.119 \\
 & Modern lossless & 939.553 & 752.336 \\
\bottomrule
\end{tabular}
\caption{Complete host times in the optimization screen. Each cell is one
observation, and all eight outputs match fresh reconstruction. The optimized
packages retain the same respective stored-state bytes.}
\label{tab:v89}
\end{table}

The measured phases explain the opposing endpoint direction without identifying
individual kernel effects. At six records, optimized compressed stage work takes
803.754 seconds versus 680.251 for modern storage. Its cheaper predecessor
validation, 5.319 versus 33.142 seconds, does not offset that difference. At twelve
records, stage work is nearly tied, 3,884.539 versus 3,891.548 seconds, while
compressed validation saves 30.030 seconds. Stage work includes several operations;
decoder time is not separately measured. Full disjoint partitions, all memory
observations and unattributed remainders are retained in E18. Including attributed
shared preparation and both requests, compressed and modern totals are
10,995.918 and 10,287.226 seconds; this attributed lifecycle also favors modern
storage. New numerical debits total 35,258 CPU seconds. All thirteen allocations
consume between 39,389.110 and 39,396.116 scheduler CPU seconds, reflecting display
precision and including both audits. The old unresolved 12,602-second reservation
is retained separately; contained clocks are not added~\evidence{E18}.

\paragraph{Untouched two-endpoint quality (V71/V72).}
The fixed teacher and both retained endpoints were evaluated on 320 articles
selected without observing losses from 2,778 eligible records in a saved bounded
frame of 3,801 complete records. The entire inspected frame, all sixty historical
quality exclusions and the separate calibration frame were protected against
reuse. A failed admission caused by incomplete short-record hash reconstruction
was preserved; a separate audit verified the exact exclusion union before the
corrected campaign. No article, model or threshold was replaced.

For each endpoint, compare the fixed-teacher candidate with sequential calibration
on the same retained records, normalization and grids. Each article contributes
127 predictions from a 128-token sequence. The six model roles share the same
320 articles. A fixed-seed paired article bootstrap uses 20,000 draws shared
across endpoints. Both pooled perplexity ratios and their one-sided 97.5\%
upper bounds must be at most 1.05; each observed maximum article ratio must be
at most 1.20. The two upper bounds use a nominal familywise level of 95\%.
This approximate bootstrap concerns the admitted frame, and the observed maximum
is not a population-tail bound.

\begin{table}[t]
\centering
\small
\begin{tabular}{rrrr}
\toprule
Retained & Pooled ratio & Upper bound & Maximum \\
\midrule
12 & 1.005421 & 1.010342 & 1.129140 \\
6 & 1.007149 & 1.011431 & 1.126961 \\
\bottomrule
\end{tabular}
\caption{Untouched perplexity ratios: fixed-teacher divided by matched sequential
calibration. All frozen guards pass. The upper bound is one-sided 97.5\%; the
maximum is over the 320 observed articles.}
\label{tab:untouched}
\end{table}

Table~\ref{tab:untouched} shows a small quality cost, not an improvement. The
paired 95\% percentile intervals are [1.000455, 1.010342] and
[1.002870, 1.011431], both above one. Candidate perplexities are 60.1053 and
60.3088, versus sequential values 59.7812 and 59.8807; the original checkpoint
and nearest-grid controls give 54.9643 and 74.6245.
All 24 CPU/CUDA reference comparisons satisfy the unchanged $10^{-8}$ mean-NLL
barrier; the maximum difference is $1.155\times10^{-14}$. The actual terminal
audit checks all 1,920 GPU and 24 CPU losses, saved sequential codes, original
parameters and recomputed statistics. A final accounting job failed because its
container lacked the scheduler client. Separate accounting recovery passed
without repeating inference, preserving that failure and its cost. Numerical
workers charge 7,526 CPU seconds; original allocations plus recovery consume
7,588.049 scheduler CPU seconds. These overlapping clocks are not additive.
The receipt and scalar-loss evidence is recorded separately~\evidence{E17}.

\paragraph{Research decision.}
The full service and optimization screen preserve a modest storage advantage
against a stronger lossless cache. Optimization reduces observed latency, but
compression still loses at one endpoint and in the attributed two-request lifecycle. Hybrid provides
a competing tradeoff when source access is allowed. These results support a
scoped space/access/latency study. Both untouched endpoints satisfy the
predeclared quality tolerances, with slightly worse perplexity than sequential
calibration. This supports the fixed target on this frame, not quality equivalence
or a general speedup. We retain the nearest-anchor failure and make no reliability
claim from single trajectory timings. Balanced repetitions would be required to
estimate a stable latency difference; they are outside the present claim.

\paragraph{Deployment limits.}
The trusted computing base includes the finite evaluator, directed arithmetic,
compiler/runtime assumptions, source provenance, and immutable inputs during
verification. Hashes are integrity commitments under that trust model, not proof
against a hostile host. The current exact evaluator runs on supported CPU
arithmetic; using a GPU feature kernel would require a separate equivalence
argument or a separately declared target. Positive ridge does not guarantee
bounded numerical certification, and refusal must remain a visible outcome.

Logical successor-state deletion leaves audit artifacts and original external
corpora intact. It does not erase information in pretrained parameters, enforce
legal data-removal obligations, or certify physical memory or filesystem erasure.
Those are separate contracts.

\section{Reproducibility and evidence status}
The review bundle provides source/protocol copies, evidence summaries E01--E18,
saved selection commitments and paired losses, typed audit-receipt derivatives,
dependency notices and a reproduction guide. E01--E15 retain the development
history and detailed costs; E16 records both disjoint-root service audits and all
ten comparisons; E17 records the actual untouched-loss and sequential-model
audits, accounting failure and separate recovery. E18 records all eight optimization
comparisons, both audits and complete host accounting. A separate optimization
source supplement extends the supplied code through E18. Its path-adapted copies
support inspection, not replay of consumed registrations. Failed attempts,
exclusions, resource ledgers and the shared eight-article exposure remain explicit.

Actual model/state/loss comparisons ran on the cluster. The portable checks verify
copied bytes, commitments and descriptive arithmetic. Generated local fixtures
mock production numerical, runtime and scheduler boundaries. Neither repeats the
empirical pipeline. Source derivation maps, raw corpus/token bodies and
model/state/runtime binaries are excluded. Independent reproduction requires
acquiring the specified external inputs and bound binaries, or separately
registering regeneration under the frozen target. The guide distinguishes these
steps from runnable offline checks; complete external-input reproduction remains
unverified. The guide records pinned model/data revisions and publisher-declared
Apache 2.0 (DistilGPT2) and ODC-BY (C4) terms, including the additional Common
Crawl conditions. These declarations do not establish rights to every underlying
web page. Upstream dependency notices accompany the package; no general public
license has been assigned to the project source.

\label{page:main-end}
\section*{Limitations}
The evidence concerns one checkpoint, short calibration records, and adaptive
development experiments. The original-checkpoint teacher was selected after the
nearest-anchor quality failure. The eight exposed articles cannot establish
quality on new data. The disjoint root tests the unchanged service. Its
untouched quality passes the frozen tolerances on one bounded article frame.
The original teacher choice remains adaptive, and prospective timing repetitions
were not performed. Two development roots do not establish broad reliability. Another checkpoint is
needed for a transfer claim.
Exact agreement concerns the declared fixed-teacher target and stored state;
it does not certify ordinary sequential calibration, GPU calibration, pretraining
unlearning, or privacy. Bounded attempts may refuse, and replay requires continued
source access. Storage savings must be assessed jointly with latency, preparation,
deployment dependencies, shared checkpoint storage, and audit costs.

\section*{Ethical considerations}
This is a research prototype for calibration-state maintenance. Exact code
recovery is not a privacy or pretraining-unlearning guarantee. Treating it as
data erasure could mislead data subjects: external corpora and audit archives
remain. Public web text may contain personal information or offensive content;
benchmark use alone does not establish its absence. We do not claim a
content-level privacy/offensive-text audit, a demographic characterization, or
fairness improvements. Raw text, token sequences and model binaries are excluded
from the review package. Independent use still requires the original input
conditions; the supplied hashes are provenance checks, not proof of anonymization.

\section*{Acknowledgements}
Generative AI assisted experiment planning, software development, analysis,
literature checking, and manuscript drafting. Human authors remain responsible
for the final methods, results, references, and text.

\begin{thebibliography}{10}
\bibitem[Frantar et~al.(2023)Frantar, Ashkboos, Hoefler, and Alistarh]{gptq}
Elias Frantar, Saleh Ashkboos, Torsten Hoefler, and Dan Alistarh.
2023. \emph{GPTQ: Accurate Post-Training Quantization for Generative Pre-trained
Transformers}. ICLR. Primary preprint:
\href{https://arxiv.org/abs/2210.17323v2}{arXiv:2210.17323, version 2}.

\bibitem[Ginart et~al.(2019)Ginart, Guan, Valiant, and Zou]{ginart}
Antonio A. Ginart, Melody Y. Guan, Gregory Valiant, and James Zou.
2019. \emph{Making AI Forget You: Data Deletion in Machine Learning}.
Advances in Neural Information Processing Systems 32.
\href{https://papers.nips.cc/paper_files/paper/2019/file/cb79f8fa58b91d3af6c9c991f63962d3-Paper.pdf}{Primary proceedings paper}.

\bibitem[Lang et~al.(2025)Lang, Helvitz, and Shlezinger]{medu}
Natalie Lang, Alon Helvitz, and Nir Shlezinger.
2025. \emph{Memory-Efficient Distributed Unlearning}.
\href{https://arxiv.org/abs/2505.03388v1}{arXiv:2505.03388, version 1}.

\bibitem[Tavory et~al.(2026)Tavory, Gade, Sarig, Touitou, and Guy]{qss}
Ami Tavory, Shripad Gade, Tal Sarig, Noam Touitou, and Ido Guy.
2026. \emph{Exact Unlearning via Quantized Sufficient Statistics}.
\href{https://arxiv.org/html/2610.07197v1}{arXiv:2610.07197, version 1}.

\bibitem[Xiong et~al.(2023)Xiong, Li, Li, and Cai]{exactfun}
Zuobin Xiong, Wei Li, Yingshu Li, and Zhipeng Cai.
2023. \emph{Exact-Fun: An Exact and Efficient Federated Unlearning Approach}.
IEEE International Conference on Data Mining, pages 1439--1444.
\href{https://doi.org/10.1109/ICDM58522.2023.00188}{Conference record};
\href{https://github.com/zuobinxiong/zuobinxiong.github.io/blob/300a0672be8d265ddebf20d6e69f790456a1bd81/assets/pdf/ExactFedUnlearning.pdf}{pinned full author version},
Theorem 2, pages 6--7. The proof comparison uses this historical author version.

\bibitem[Zhang et~al.(2023)Zhang, Song, and Sun]{qebverif}
Yedi Zhang, Fu Song, and Jun Sun.
2023. \emph{QEBVerif: Quantization Error Bound Verification of Neural Networks}.
\href{https://arxiv.org/abs/2212.02781v2}{arXiv:2212.02781, version 2}.

\bibitem[Zhao et~al.(2026)Zhao, Wang, Chang, and He]{execcert}
Ziyu Zhao, Xinyu Wang, Xiao-Wen Chang, and Yixuan He.
2026. \emph{From Mathematical to Executable Certificates for Machine Unlearning}.
\href{https://arxiv.org/html/2610.02268v1}{arXiv:2610.02268, version 1}.

\bibitem[Hugging Face(n.d.)]{distilgpt2}
Hugging Face. n.d. \emph{DistilGPT2 model card}.
Pinned model repository revision 2290a62682d06624634c1f46a6ad5be0f47f38aa.
\href{https://huggingface.co/distilbert/distilgpt2/blob/2290a62682d06624634c1f46a6ad5be0f47f38aa/README.md}{Publisher model card}.

\bibitem[Raffel et~al.(2020)Raffel, Shazeer, Roberts, Lee, Narang, Matena, Zhou, Li, and Liu]{c4}
Colin Raffel, Noam Shazeer, Adam Roberts, Katherine Lee, Sharan Narang,
Michael Matena, Yanqi Zhou, Wei Li, and Peter J. Liu.
2020. \emph{Exploring the Limits of Transfer Learning with a Unified Text-to-Text
Transformer}. Journal of Machine Learning Research, 21(140):1--67.
\href{https://jmlr.org/papers/v21/20-074.html}{Primary journal article}.

\bibitem[Merity et~al.(2016)Merity, Xiong, Bradbury, and Socher]{wikitext}
Stephen Merity, Caiming Xiong, James Bradbury, and Richard Socher.
2016. \emph{Pointer Sentinel Mixture Models}.
\href{https://arxiv.org/abs/1609.07843v1}{arXiv:1609.07843, version 1}.

\end{thebibliography}
\typeout{SECOND-MAIN-END-PAGE=\getpagerefnumber{page:main-end}}
\end{document}
